\chapter{Conclusion}

The objectives set out in the introduction have been successfully achieved. We found a way to create a videoconferencing node based entirely on freely available components. The tools developed during this work make it possible to prepare and record a conference. The way in which we combined the existing components is unique within the Access Grid environment.

\paragraph*{We consider the following to be the main contributions:}
\begin{enumerate}
%---------------------------
\item Creation of a videoconferencing node at the Ciant organization. It is the second Access Grid node in the Czech Republic and can serve as a model for nodes at other organizations.

%---------------------------
\item As part of creating the node, a website was developed. In addition to information about the node itself, it contains a collection of tools related to videoconferencing technologies. The engine used to create the website is available to interested users.

%---------------------------
\item The contacts established through the node's activities to date will be used in future cooperation between MFF UK and the MARCEL network. This cooperation could generate further interesting projects that could serve as topics for bachelor's, master's, and doctoral theses.

%---------------------------
\item The practical experience gained while creating the node made it possible to design an Access Grid node composed entirely of freely available components. Applications that are otherwise available on commercial platforms had to be provided separately.

%---------------------------
\item Creation of a shared multimedia player for the Linux operating system—SharedMPlayer—which had previously been missing from Access Grid. The repair of the existing version of SharedPresentation is also worth mentioning.

%---------------------------
\item Another software component developed as a result of this work is EVIC, a modified version of the VIC application from the VRVS project. The main objective was to find a suitable candidate for recording a VNC screen, which the original VIC developed at UCL did not support. EVIC became part of the StreamVNC tool and provides control elements located on a remote X11 display.

%---------------------------
\item The integrated StreamVNC tool was developed to prepare presentations suitable for broadcasting to an Access Grid conference. It enables local recording of VNC content, audio, and video from a camera. The data stream, in the form of RTP packets, is stored using the AGVCR tool.

%---------------------------
\item From a theoretical perspective, we described the main problems of Access Grid and videoconferencing technologies in general. Using small examples, we demonstrated ways to traverse NAT and firewalls, as well as tunnelling and multicast-network bridging. For the node at Ciant, we built QuickBridge, which partially addresses the problem of multicast-network unavailability. In the future, we are considering deploying the newer version AG 3, which provides improved bridging using UMTP.

\end{enumerate}

A substantial part of these results may lead to new projects in the future. Videoconferencing technologies are constantly evolving, and users continue to make new demands regarding transmission quality, functionality, and ease of use. Access Grid needs new, high-quality audio and video codecs, which implies improvements to the existing VIC and RAT components. Following the InSORS and VRVS models, it would be appropriate to implement new bridges (SIP, Skype, H.323, Asterisk \ldots). We also propose adapting the current robust architecture for use in constrained environments without negatively affecting transmission quality at other points in the network. Usability could be improved by simplifying the VenueClient interface, for example, following the EVO, VRVS, or InSORS models.

\medskip
We hope that the Ciant AG node project will continue, and we welcome interested users to visit our wiki pages \cite{CIANT-AGnodeWiki}, where the latest versions of all the components mentioned above can be found, as well as this text in electronic form.
\proofed
